% =====================================================================
%  ch3d_hamiltonian.tex — §3.4 无穷 Hamilton 扰动；§3.5 continuation
% =====================================================================
\subsection{无穷 Hamilton 扰动}\label{ss:FMH}

现在描述 $M$ 的 Fukaya 范畴的一个版本，它可以用局部片段来表达。这一构造涉及（在某种意义下）很大的——「无穷」——Hamilton 扰动，与 H.~Lee 的学位论文 \cite{Lee} 非常相似。Lee 把黎曼曲面分解成裤子，而我们改用顶点的邻域（即裤子的\emph{镜像}）作为构件。

对每条半边 $e/v\in E$，选取 $M_e$ 与 $[0,4]\times S^1/\!\sim$ 的等同，其中 $\{0\}\times S^1$（相应地 $\{4\}\times S^1$）等同于 $v$（相应地 $v'$），使得辛形式为
\[
\omega=\frac{A_e}{4}\,d\tau\wedge d\psi,
\]
这里 $\tau$ 与 $\psi$ 分别是 $[0,4]$ 与 $S^1=\R/\Z$ 上的坐标。

我们假定在 $\tau=1$ 附近与 $\tau=3$ 附近（事实上，在下面引入的进一步扰动的整个支集上），定义范畴 $\cF(M)$ 时所用的 Hamilton 函数 $h$ 可以写成 $\tau$ 坐标的函数。取一列 $C^\infty$ 函数 $f_n:[0,4]\to\R$，它们在 $\tau=1$ 与 $\tau=3$ 之外为常数，并收敛到连续函数 $f:[0,4]\to\R$，使得：
\begin{itemize}
\item[(i)] $f=f_n=0$ 在 $\tau=0$ 与 $\tau=4$ 附近，且 $f$ 与 $f_n$ 在 $\tau=2$ 附近为常数；在 $[0,1-\frac1n]\cup[3+\frac1n,4]$ 上 $f_n=f$，而在 $[1+\frac1n,3-\frac1n]$ 上 $f_n-f$ 为常数；
\item[(ii)] 在 $[0,1)\cup(3,4]$ 上 $f''_n\leq 0$，在 $(1,3)$ 上 $f''_n\geq 0$（从而 $f''$ 也有同样性质）；
\item[(iii)] $f'_n(1)=-n$，$f'_n(3)=n$，且 $\lim_{\tau\to 1}f'(\tau)=-\infty$，$\lim_{\tau\to 3}f'(\tau)=+\infty$。
\end{itemize}

我们考虑 Hamilton 扰动
\[
H_n=\varepsilon h+\frac{A_e}{4}f_n(\tau).
\]
$h$ 在 $f'_n$ 的支集上只依赖于 $\tau$ 这一假设，保证 $h$ 与 $f_n$ 生成的 Hamilton 流可交换，并且 $H_n$ 的 time-$1$ 流与 $\varepsilon h$ 的 time-$1$ 流相差一个旋转 $\psi\mapsto\psi+f'_n(\tau)$。范畴 $\cF(M;H)$ 用 $H_n$ 代替 $\varepsilon h$ 作为 Floer 复形的 Hamilton 扰动来定义，并按下面讨论的方式取 $n\to\infty$。

我们对 $\cF(M;H)$ 的对象再加一些条件。对 v.b. 型拉格朗日，假定坐标 $\tau$ 在每个分量上严格单调（从而拉格朗日只经过 $\tau=1$ 与 $\tau=3$ 处的「颈」一次）；注意 $\cF(M)$ 的每个 v.b. 型对象都同构于满足此条件的对象。还假定 $CF^*(L,L';\varepsilon h)$ 的生成元都避开 $\tau=1$ 与 $\tau=3$。点型拉格朗日对我们的论证并非必需，但只要它们与 $\tau=1$、$\tau=3$ 处的圆不相交就可以允许；这就排除了支撑在我们分解各片边界上的对象。

给定一对 v.b. 型对象 $(L,\cE),(L',\cE')$，考虑 Floer 复形 $CF^*((L,\cE),(L',\cE');H_n)$，其生成元以 $H_n$ 的 Hamilton 向量场从 $L$ 出发、终止于 $L'$ 的 time-$1$ 轨道之集 $\cX(L,L';H_n)$ 为指标，等价地，以 $\phi^1_{H_n}(L)$ 与 $L'$ 的交点为指标。对任何 $n$，都可以在 §\ref{s:Amodel} 的构造中用 $H_n$ 代替 $\varepsilon h$，得到与 $\cF(M)$ 拟等价的 $A_\infty$-范畴 $\cF(M;H_n)$。然而，由于对给定输入下传播圆盘的次数没有\emph{先验}界，H.~Lee 的论证 \cite{Lee} 不允许我们断定：对任何有限的 $n$，$A_\infty$-操作 $\mu^k_{H_n}$ 都能由局部考虑理解，即使只限于有限个对象。

为处理这一点，把 $CF^*((L,\cE),(L',\cE');H)$ 定义为如下可数无穷维 $K$-向量空间的完备化：其生成元对应于 $H=\varepsilon h+\frac{A_e}{4}f(\tau)$ 的 Hamilton 向量场在流未定义的圆 $\tau=1$ 与 $\tau=3$ 之外的补集中的 time-$1$ 轨道（以及局部系统在轨道端点处纤维之间的态射）。也就是说，$CF^*(L,L';H)$ 由形如 $\sum \rho_p$ 的形式和组成，其中 $\rho_p\in\Hom(\cE_{p},\cE'_{p'})$ 且 $\|\rho_p\|\to 0$（即 $\val(\rho_p)\to+\infty$）。

Floer 复形 $CF^*((L,\cE),(L',\cE');H_n)$ 在 $n\to\infty$ 时以下述意义稳定。对 $\delta>0$，令
\[
\cN_{\delta}=\bigcup_{e\in E(G)}\cN_{e,\delta},\qquad
\cN_{e,\delta}=([1-\delta,1+\delta]\cup[3-\delta,3+\delta])\times S^1\subset M_e .
\]
则 $CF^*((L,\cE),(L',\cE');H_n)$ 中位于 $\cN_{1/N}$ 之外的生成元对一切 $n\geq N$ 完全相同，于是可以把 $CF^*((L,\cE),(L',\cE');H)$ 看作这些 Floer 复形的朴素极限的完备化。（由于我们的圆盘计数按辛面积而非拓扑能量加权，对大的 $n$ 可以直接把 Floer 生成元等同起来，无需 \cite{AuSm} 中讨论的作用量重新标定。）

考虑扰动所诱导的旋转 $\psi\mapsto\psi+f'_n(\tau)$ 的效果，可以看到：在 $L$ 与 $L'$ 于 $\tau=1$、$\tau=3$ 附近的几何的温和假设下，集合 $\cX(L,L';H_n)$（相应地 $\cX(L,L';H)$）与 $\cX(L,L';\varepsilon h)$ 的差别在于添加了：
\begin{itemize}
\item $(0,1)\times S^1$ 中 $n$ 个（相应地，无穷多个）次数 $1$ 的生成元；
\item $(1,3)\times S^1$ 中 $2n$ 个（相应地，「两倍」无穷多个）次数 $0$ 的生成元；
\item $(3,4)\times S^1$ 中 $n$ 个（相应地，无穷多个）次数 $1$ 的生成元。
\end{itemize}

固定 v.b. 型拉格朗日 $L_0,\dots,L_k$ 与输入生成元 $p_i\in\cX(L_{i-1},L_i;H_n)$。考虑对 Hamilton 函数 $H_n$ 的传播扰动全纯圆盘中映到 $M_e$ 的一个分量 $u_e:D_e\to M_e$，并设 $([0,1)\times S^1/\!\sim)\subset M_e$ 含有 $u_e$ 的像的一部分，但不含其输出。则 $u_e$ 到 $M_e-\{p_v,p_{v'}\}$ 的万有覆盖的提升沿 $\psi$ 方向有一个最大「宽度」，它由 $u_e$ 的输入决定，并且对那些映到 $\tau=0$ 处节点 $p_v$ 的输入而言，还由 $u_e$ 在节点附近带状端中的局部次数决定。然而扰动 $H_n$ 阻止 $u_e$ 任何宽度小于 $n$ 的部分从输入到输出沿 $\tau$ 增大的方向穿过 $\tau=1$。因此，若 $n$ 相对于 $u_e$ 在映到 $p_v$ 的输入处的局部次数之和足够大，就得到矛盾，$u_e$ 的输出也必须位于 $\tau<1$ 处；见 \cite[Lemma 3.5]{Lee}（并见 \cite[Proposition 5.5]{AuSm}）。

于是得到如下结论。对 $G$ 的每个顶点 $v$，记 $P_v$ 为各半边 $e/v$ 对应的子集 $([0,3]\times S^1/\!\sim)\subset M_e$ 之并。对每条边 $e$，记 $N_e\subset M_e\subset M$ 为子集 $[1,3]\times S^1$。（于是当 $e$ 是连接 $v$ 与 $v'$ 的唯一一条边时，$N_e=P_v\cap P_{v'}$。）则：

\begin{prop}\label{prop:hlee}
给定对 Hamilton 函数 $H_n$ 的任一传播扰动全纯圆盘 $u:S\to M$，其边界在 $L_0,\dots,L_k$ 上，输入为 $p_i\in\cX(L_{i-1},L_i;H_n)$，则下列之一成立：
\begin{itemize}
\item $u$ 的像完全含于某个 $P_v$（$v\in V(G)$）之内，且输出标记点位于所有 $N_e$（$e/v$）之外；
\item $u$ 的像完全含于某个 $N_e$（$e\in E(G)$）之内；
\item 至少有一个输入生成元 $p_i$ 落在 $\cN_{k/n}$ 内；
\item 圆盘 $u$ 穿过 $M$ 的某个节点传播，其输出次数 $k_{out}\geq n/k$。
\end{itemize}
\end{prop}

最后一种情形中，输出次数 $\geq n/k$ 的传播意味着圆盘的辛面积有常数倍 $n$ 的下界。因此有：

\begin{prop}\label{prop:lowarealocal}
对给定的一族输入生成元 $p_i\in\cX(L_{i-1},L_i;H)$ 与常数 $A>0$，存在 $N=N(A)$，使得当 $n\geq N$ 时，任何以 $p_1,\dots,p_k$ 为输入且面积 $\leq A$ 的传播扰动全纯圆盘都完全落在单个片段 $P_v$（或 $N_e$）之内，且其输出位于 $\cN_{1/N}$ 之外。此外，这类圆盘的模空间对一切 $n\geq N$ 彼此双射。
\end{prop}

\begin{proof}
命题的第一部分由命题~\ref{prop:hlee} 立得，因为当 $n$ 足够大时，面积界排除了输出次数 $\geq n/k$ 的传播。此外，传播次数的界给出传播圆盘各分量沿 $\psi$ 坐标的「宽度」的界，从而也给出输出的界；而 $\tau=1$、$\tau=3$ 附近的生成元对应于 Hamilton 流沿 $S^1$ 方向缠绕越来越多的轨道。最后，对不穿过 $\tau=1$ 的圆盘，不同 $n\geq N$ 时传播圆盘模空间之间的双射是显然的；对穿过 $\tau=1$ 的圆盘（必沿 $\tau$ 减小的方向从输入到输出），把扰动 Cauchy--Riemann 方程的解改写为边界在 $L_i$ 经 Hamilton 流作用的像上的多边形之后，就清楚地看到：增大 $n$ 只是通过加宽穿过 $\tau=1$ 处颈部的带状部分来变形这些多边形。（相关论证并见 \cite[Section 3]{Lee} 与 \cite[Section 5]{AuSm}。）
\end{proof}

这使我们能把 $\F(M;H)$ 中的 $A_\infty$-操作定义为用 Hamilton 函数 $H_n$ 的操作的朴素极限：给定 $p_i\in\cX(L_{i-1},L_i;H)$ 与酉元 $\rho_i\in\hom(\cE_{i-1|p_i},\cE_{i|p'_i})$，定义
\begin{equation}\label{eq:mu_k_H}
\mu^k_H(\rho_k,\dots,\rho_1)=\lim_{n\to \infty} \mu^k_{H_n}(\rho_k,\dots,\rho_1),
\end{equation}
即 $CF^*((L_0,\cE_0),(L_k,\cE_k);H)$ 中与 $\mu^k_{H_n}(\rho_k,\dots,\rho_1)$ 对一切 $n>N(A)$ 模 $T^A$ 相同的元素。再按线性性把它延拓到生成元的有限和，然后按连续性延拓到完备化态射空间中的任意输入。

具体地说，$\mu^k_H(\rho_k,\dots,\rho_1)$ 可以理解为传播圆盘的带权计数，其中 Hamilton 扰动相对于给定输入以及 $S$ 各节点处的局部次数 $k_{out}$ 取得足够大；由命题~\ref{prop:hlee}，这些圆盘都落在单个 $P_v$ 之内，因而只能穿过 $M$ 的一个节点传播。

有一点小技术性说明：在上述构造中，我们对 $\mu^k_{H_n}$ 的输入与输出使用同一个 Hamilton 函数 $H_n$，这意味着当 $k\geq 2$ 时，扰动 Cauchy--Riemann 方程涉及一个不闭的 $1$-形式 $\beta$（对紧 $M$ 这不成问题，因为 $H_n$ 有界；在缠绕情形则应对 $M$ 的非紧分量改用 Abouzaid 的重新标定技巧）。不过也可以像 \cite{Lee} 与 \cite{AuSm} 那样，在 Cauchy--Riemann 方程中使用闭 $1$-形式，并让 $\mu^k$ 取值于 Hamilton 扰动为 $kH_n$ 的 Floer 复形，其几何行为与 $H_{kn}$ 本质相同。此时 $n\to\infty$ 极限构造的细节有所不同（如果对不同 Hamilton 函数的复形的链层极限引入「望远镜」模型，可能更为复杂），但即便如此，在对拉格朗日 $L_i$ 的温和几何假设下，仍可把构造改写成（完备化的）朴素极限。

我们指出命题~\ref{prop:hlee} 的如下推论，它在 §\ref{s:equivalence} 中要用到：

\begin{prop}\label{prop:ainfty_ideal}
对每个 $v\in V(G)$，Floer 复形中位于 $P_v$ 之外的生成元所（完备化地）张成的子空间构成 $\cF(M;H)$ 中的一个 $A_\infty$-理想。记 $\cF(P_v;H)$ 为 $\cF(M;H)$ 对此 $A_\infty$-理想的商。类似地，对每条边 $e$，位于 $N_e$ 之外的生成元张成的子空间构成 $\cF(M;H)$ 中的 $A_\infty$-理想，记 $\cF(N_e;H)$ 为相应的商。
\end{prop}

\subsection{Continuation $A_\infty$-同态}\label{ss:continuation}

本节最后用拉格朗日 Floer 理论中的 continuation 映射（例如见 \cite{SeBook}）构造从 $\cF(M)$ 到 $\cF(M;H)$ 的 $A_\infty$-同态。由于我们的比较论证依赖另一条途径（见 §\ref{ss:HPLFloer}），这里略去构造高阶项时的一些细节。

我们构造 $A_\infty$-同态 $\mathfrak{K}_n:\cF(M)\to \cF(M;H_n)$，其 $k$ 阶项
\[
\mathfrak{K}_n^k:\bigotimes_{i=1}^k CF^*((L_{i-1},\cE_{i-1}),(L_i,\cE_i);\varepsilon h)\to
CF^{*+1-k}((L_0,\cE_0),(L_k,\cE_k);H_n)
\]
计数带 $k$ 个输入的扰动传播全纯圆盘，其中的 Hamilton 扰动在输入端为 $\varepsilon h$、在输出端为 $H_n$。

一阶映射 $\mathfrak{K}_n^1$ 最容易描述。固定一簇光滑 Hamilton 函数 $H_\sigma$（$\sigma\in\R_{\geq 0}$），使得 $\sigma=0$ 时 $H_\sigma=\varepsilon h$，$\sigma=n$ 时 $H_\sigma=H_n$；再固定光滑非增函数 $\sigma:\R\to\R$，使得在 $(-\infty,-1)$ 上 $\sigma=n$、在 $(1,\infty)$ 上 $\sigma=0$。带单个输入的传播圆盘的定义域是一列每个带两个标记点的圆盘（即带状域 $\R\times[0,1]$）构成的链，$S=D_1\cup\dots\cup D_\ell$（$D_1$ 承载输入标记点 $z_1$，$D_\ell$ 承载输出 $z_0$）。我们考虑映射 $u:S\to M$，其中一个分量 $D_j$ 满足通常的 Floer continuation 方程
\[
(du-X_{H_{\sigma(s)}}\,dt)^{0,1}=0,
\]
其在输入端（$s\to+\infty$）的 Hamilton 函数为 $\varepsilon h$、在输出端（$s\to-\infty$）为 $H_n$；而沿输入到输出方向位于其前（相应地，后）的分量 $D_1,\dots,D_{j-1}$（相应地，$D_{j+1},\dots,D_\ell$）是 Hamilton 函数 $\varepsilon h$（相应地，$H_n$）的扰动全纯带状域。对这类刚性（即属于期望维数为零的解模空间，等价地，输入与输出生成元次数相同）的扰动传播圆盘（对 continuation 发生所在分量 $S$ 的所有可能选择）作带符号与权的计数，符号与权如同 $A_\infty$-操作的定义，就得到映射
\[
\mathfrak{K}_n^1:CF^*((L_0,\cE_0),(L_1,\cE_1);\varepsilon h)\to
CF^*((L_0,\cE_0),(L_1,\cE_1);H_n),
\]
通过考察一维模空间容易验证它是链映射。

\begin{remark}\label{rmk:continuation}
尽管定义允许 Hamilton 函数在传播圆盘的任一分量上改变，但正则刚性 continuation 轨道的各分量本身是刚性的；这意味着：若输入与输出分别是相应 Floer 复形的次数 $1$、次数 $0$ 生成元，则 continuation 实际上必须发生在输入分量 $D_1$（相应地，输出分量 $D_\ell$）上。

事实上，在我们的设定中，从 $M_e$ 内部的次数 $1$ 生成元出发的 continuation 轨道必为常值。因此 $\mathfrak{K}_n^1$ 在 $CF^1$ 上就是朴素包含；而对次数 $0$ 生成元，它与朴素包含（常值轨道）相差一些扰动传播全纯带状域的计数，其中输出分量是通常意义下从 $\varepsilon h$ 到 $H_n$ 的 continuation 轨道，其余分量是 Hamilton 函数 $\varepsilon h$ 的扰动全纯带状域。
\end{remark}

此外，与上一节相同的论证表明：辛面积（或能量——两者等价，因为 $f_n$ 与 $f$ 一致有界）有界的 continuation 轨道，从而通过 $M$ 各节点时传播次数有界的轨道，在 $n\to\infty$ 时必须稳定，即对一切足够大的 $n$，模空间相同。这使我们能定义
\[
\mathfrak{K}^1:CF^*((L_0,\cE_0),(L_1,\cE_1);\varepsilon h)\to
CF^*((L_0,\cE_0),(L_1,\cE_1);H),\qquad
\mathfrak{K}^1(\rho)=\lim_{n\to\infty}\mathfrak{K}^1_n(\rho).
\]
在恒等式 $\mu^1_{H_n}\circ\mathfrak{K}^1_n=\mathfrak{K}^1_n\circ\mu^1$ 中取 $n\to\infty$，即知 $\mathfrak{K}^1$ 也是链映射。

由注~\ref{rmk:continuation} 可知：对次数 $0$ 的 Floer 生成元，$\mathfrak{K}^1$ 计数的是这样的传播圆盘，其中 $S$ 的输出分量是从 $\varepsilon h$ 到 $H$（即对足够大的 $n$，到 $H_n$）的 continuation 轨道，其余分量是 Hamilton 函数 $\varepsilon h$ 的 Floer 方程的解；而对次数 $1$ 的生成元，continuation 发生在输入端。

$A_\infty$-同态 $\mathfrak{K}_n$ 的高阶项涉及：对每个稳定结点定义域 $S=\bigsqcup D_j/\!\sim$（并在这些定义域的模空间上连续且一致地），选取一参数族的 Hamilton 扰动数据，使得在该族的一端，除输出标记点 $z_0$ 附近的带状端（那里发生到 $H_n$ 的 continuation）之外，Hamilton 函数处处为 $\varepsilon h$；而在该族的另一端，除输入标记点 $z_1,\dots,z_k$ 附近的带状端之外，Hamilton 函数处处为 $H_n$。实现这一点的一种方式是：对每个 $S$ 选取光滑函数 $s:S-\{z_0,\dots,z_k\}\to\R$，使得在输出标记点处 $\lim_{z\to z_0}s(z)=-\infty$，在输入标记点处 $\lim_{z\to z_i}s(z)=+\infty$，并且在 $S$ 的每个分量上，$s$ 从输入到输出单调递减。这个选取应当在稳定结点圆盘的模空间上连续，并对定义域的退化保持一致。然后考虑含 Hamilton 函数 $H_{\sigma(s(z)-s_0)}$ 的 Floer continuation 方程的解，其中参数 $s_0\in\R$ 可自由变化，决定 Hamilton 扰动从 $\varepsilon h$ 切换到 $H_n$ 处的 $s$ 的水准集。

由于 $H_n=\varepsilon h$ 在 $M$ 的节点附近成立（而且可以保证所有 $H_\sigma$ 也如此），当 $s_0$ 扫过 $S$ 某节点处的 $s$ 值时，continuation 扰动行为的细节并不特别重要。需要更小心的是 $S$ 的某些分量为不稳定（带状）的情形：若 continuation 同时跨越 $S$ 的若干个不稳定分量，最显然的构造无法计入定义域的自同构。概念上最简单的办法是给 $S$ 的每个不稳定分量添加一个边界标记点来稳定定义域，并要求相应分量 $M_e$ 的 $\tau$ 坐标取指定值。（或者，像注~\ref{rmk:continuation} 那样考察刚性 continuation 构形的结构，可以排除若干可能情形，并对剩余情形给出特别定义。）

与线性项的情形一样，注意到面积低于固定阈值的传播圆盘对 $\mathfrak{K}_n^k$ 的贡献对足够大的 $n$ 稳定，便可取 $n\to\infty$ 的极限，令
\[
\mathfrak{K}^k(\rho_k,\dots,\rho_1)=\lim_{n\to \infty} \mathfrak{K}^k_n(\rho_k,\dots,\rho_1).
\]

我们断言 $A_\infty$-函子 $\mathfrak{K}:\cF(M)\to\cF(M;H)$ 是拟等价。确立此类结论的通常方法是：把 $H$ 与 $\varepsilon h$ 的角色对调，用 Floer 理论中的 continuation 映射构造反方向的另一个 $A_\infty$-函子，并通过同伦论证说明它是 $\mathfrak{K}$ 的拟逆。我们预期这在本设定中也可行，但更简便的做法是：只需说明 $A_\infty$-函子 $\mathfrak{K}$ 的线性项是链复形的拟同构；这将由 §\ref{ss:HPLFloer} 的论证给出——那里将说明 $\mathfrak{K}^1$ 与一个基于同调扰动引理的纯代数构造一致。
% [Tutor Note: 点拨] §3.4 的全部努力只为一件事：让 $\mu^k$ 的贡献「局部化」到单个 $P_v$（顶点处的裤子镜像）或单个颈 $N_e$ 上。方法是把 Hamilton 函数 $H_n$ 在 $\tau=1,3$ 处的斜率取到 $-n,+n$，使得任何宽度小于 $n$ 的圆盘无法穿过颈部；而输出次数 $k_{out}$ 又给出面积下界。于是对固定面积截断、$n$ 足够大时，圆盘只能落在局部片段里——这正是 §5 用 Čech 型局部化证明等价的基础。依据：原文 Prop 3.16–3.18。
